\begin{center}
\textsc{Gauss} an \textsc{Gerling}.\quad Göttingen, 19.~Juli 1827.
\end{center}

\par\noindent\hspace*{\parindent}\mbox{.\ .\ .\ .\ .\ .}\quad Obwohl das Wetter meine astronomischen Arbeiten nicht in
dem Maße begünstigt hat, wie ich es mir gewünscht hatte, glaube ich
doch, die Amplitude des Bogens zwischen den Sternwarten in Göttingen
und Altona bis auf einen sehr kleinen Bruchteil einer Bogensekunde bestimmt zu
haben. Ich habe sehr viele Sterne herangezogen und
insgesamt gegen 900 Beobachtungen gemacht. Es bleibt der Unterschied von
\par\clearpage
\noindent\null\hfill{\small\textbf{ZUR HANNOVERSCHEN TRIANGULATION.}}\hfill\makebox[0pt][r]{\fontsize{13}{16}\selectfont 391}\par\medskip
\noindent \(5^{\prime\prime}\) gegenüber der Berechnung aus den Dreiecken mit \textsc{Walbecks} Erddimensionen,
und es wäre daher sehr interessant gewesen, noch einen Zwischenpunkt zu
haben. Mit dem Sektor war das unter den gegebenen Umständen unmöglich.
Vielleicht beobachten wir aber künftig noch gemeinsam in Celle mit einem im ersten Vertikal
aufzustellenden Passageninstrument, ebenso wie \textsc{Schumacher}
vorher in Altona und ich nachher in \mbox{Göttingen.\quad\mbox{.\ .\ .\ .\ .\ .}}\quad Noch vor meiner
Reise habe ich auch sämtliche bayerischen Dreiecke mitgeteilt
bekommen, die durch Ihre mit meinen verknüpft \mbox{sind.\quad\mbox{.\ .\ .\ .\ .\ .}}
\par
